% -------------------------------------------------------------- %
%      csd5748 - HY-225 - Spring 2025 - CSD-UOC - Ex. 08-09      %
% -------------------------------------------------------------- %

\documentclass{article}

\usepackage[utf8]{inputenc} 
\usepackage[english, greek]{babel}
\usepackage{subcaption}
\usepackage{hyperref}
\usepackage{xcolor}
\hypersetup{ hidelinks, colorlinks=true, linkcolor=blue, urlcolor=blue }
\captionsetup[subfigure]{labelformat=empty}
\usepackage{amsmath}
\usepackage{amsfonts}
\usepackage{amssymb}
\usepackage{array}
\usepackage{graphicx}
\usepackage[a4paper, left=1.1cm, right=1.1cm, top=2cm, bottom=2cm]{geometry}
\usepackage[absolute, overlay]{textpos}
\usepackage{placeins}
\usepackage{listings}
\usepackage{xcolor}
\setlength{\jot}{0pt}
\setlength{\mathsurround}{0.1cm}
\setlength{\jot}{0.1cm}
\newcommand{\skato}{\hspace{0.25cm}}
\newcommand{\skata}{\hspace{1cm}}
\usepackage[english, greek]{babel}

\begin{document}

% -------------------------------------------------------------- %

\text{   }
\vspace{-1.0cm}
\begin{center}
\setlength{\fboxrule}{2pt}
\setlength{\fboxsep}{15pt} 
\fbox
{
      {\noindent {\Huge \textbf{\textlatin{HY225 -} \textbf{Ασκήσεις 12} \textlatin{- csd5748}}}}
}
\end{center}
\vspace{1cm}


% -------------------------------------------------------------- %


\begin{center}
    {\underline{\Large \textbf{Άσκηση 12.4}}}
\end{center}
\vspace{0.12345cm}



\subsubsection*{\textbf{Δεδομένα:}}
\begin{itemize}
    \item Εικονικές διευθύνσεις: 20 \textlatin{bits} (1\textlatin{MB} χώρος διευθύνσεων ανά διεργασία).
    \item Μέγεθος σελίδας: 4\textlatin{KB} ($2^{12}$ \textlatin{bytes}) άρα 12 (\textlatin{LS bits} κάθε εικονικής διεύθυνσης για \textlatin{offset}.
    \item Αριθμός εικονικών σελίδων: 1 \textlatin{MByte} / 4 \textlatin{KBytes} = 256 \ σελίδες για κάθε διεργασία. 
    \item Φυσική μνήμη: 64\textlatin{KB} (16 \textlatin{bits} φυσικές διευθύνσεις).
    \item Αριθμός φυσικών σελίδων: 64\textlatin{KB} / 4\textlatin{KB} = 16 \ σελίδες.
    \item Αριθμός εικονικής σελίδας: 20 - 12 = 8 \textlatin{MS bits} χρησιμοποιούνται
σαν \textlatin{index} στον πίνακα μετάφρασης της διεργασίας.
    \item Θέσεις πίνακα μετάφρασης: \ \textlatin{Index 8bits}, \ \textlatin{Physical Page 4bits}, \ \textlatin{Valid bit}, \ \textlatin{Prot. 3bits}, \ \textlatin{Dirty bit}, \ \textlatin{Ref. bit}.
    \item Εικονική διεύθυνση: π.χ. \textlatin{FE210}.
\end{itemize}


\subsubsection*{\textbf{α)}}
\subsubsection*{\textbf{Ζητούμενα:}}
\begin{itemize}
    \item Καταχωρητή που περιέχει τον \textlatin{pointer} στον πίνακα μετάφρασης της παρούσας
διεργασίας.
    \item Τον πίνακα μετάφρασης.
    \item Την εικονική διεύθυνση (20 \textlatin{bits}) που γεννά ο επεξεργαστής και τα πεδία από τα οποία αποτελείται.
    \item Από που προέρχεται το \textlatin{index} στον πίνακα μετάφρασης και τι διαβάζουμε από τη θέση εκείνη του πίνακα.
    \item Πώς συνθέτουμε τη φυσική διεύθυνση (16 \textlatin{bits}).
\end{itemize}


\vspace{0.345cm}

\[
    \boxed{ \ \textlatin{Register: \ *ptr \ (pointer to the beginning of the translation table below ...)} } \hspace{4.555cm}
\]

\vspace{0.25cm}

\[
256 \times 6 \ 
\left\{ \ \
\begin{aligned}
    &\smash{\raisebox{2.5em}{\textlatin{*ptr \ \rightarrow \ \ \ }}} 
\begin{tabular}{|c|c|c|c|c|c|}
    \hline
     \textbf{\textlatin{Index}} & \textbf{\textlatin{Physical Page}} & \textbf{\textlatin{Valid Bit}} & \textbf{\textlatin{Protection bits}} & \textbf{\textlatin{Dirty bit}} & \textbf{\textlatin{Reference bit}}  \\ \hline
     00 & \textlatin{Null} & \textlatin{x} & \textlatin{xxx} & \textlatin{x} & \textlatin{x}  \\ \hline
     01 & \textlatin{Null} & \textlatin{x} & \textlatin{xxx} & \textlatin{x} & \textlatin{x}  \\ \hline
     \cdots & \cdots & \cdots & \cdots & \cdots & \cdots  \\ \hline
     \textlatin{FE} & \textlatin{A} & 1 & \textlatin{xxx} & \textlatin{x} & \textlatin{x}  \\ \hline
     \cdots & \cdots & \cdots & \cdots & \cdots & \cdots  \\ \hline
     \textlatin{FF} & \textlatin{Null} & \textlatin{x} & \textlatin{xxx} & \textlatin{x} & \textlatin{x}  \\ \hline
\end{tabular}_{\textlatin{\ TT}}
\end{aligned}
\right.
\]

\vspace{0.35cm}

\[
    \boxed{ \ \textlatin{Virtual Address: \ 20 bits \ = \ 8MS \ \& \ 12LS \ = \ Page Number \ \& \ Ofset Byte \ = \ FE \ \& \ 210} \ }
\]

\vspace{0.5cm}

\text{- Το \textlatin{index} στον πίνακα μετάφρασης προέρχεται από τα 8 \textlatin{MS bits} της εικονικής διεύθυνσης.}

\text{- Από εκείνη την θέση του πίνακα μετάφρασης διαβάζουμε τη σελίδα που αντιστοιχεί στην
φυσική σελίδα.}

\text{- Η φυσική διεύθυνση προκύπτει από τα 12 \textlatin{LS bits} της εικονικής διεύθυνσης
και από τα 4 \textlatin{bits} της φυσικής σελίδας.}

\ \ (\text{A210}) \

% -------------------------------------------------------------- %



\newpage


\subsubsection*{\textbf{β)}}
 
\subsubsection*{\textbf{Ζητούμενο:}}
\begin{itemize}
    \item Δείξτε τα περιεχόμενα του πίνακα μετάφρασης της διεργασίας μας (256 γραμμές επί \underline{4} πεδία ανά γραμμή) \ \textlatin{???}
\end{itemize}

\vspace{0.5cm}

\[
\begin{tabular}{|c|c|c|c|c|}
    \hline
     \textbf{\textlatin{Index}} & \textbf{\textlatin{Physical Page}} & \textbf{\textlatin{Valid Bit}} & \textbf{\textlatin{Protection bits}} & \textbf{\textlatin{Dirty bit}}  \\ \hline
     00 & \textlatin{Null} & 0 & 000 & \textlatin{x}  \\ \hline
     01 & 4 & 1 & 101 & \textlatin{x}  \\ \hline
     02 & \textlatin{Null} & 0 & 101 & \textlatin{x} \\ \hline
     \cdots & \textlatin{Null} & 0 & 101 & \textlatin{x} \\ \hline
     09 & \textlatin{Null} & 0 & 101 & \textlatin{x} \\ \hline
     \textlatin{0A} & \textlatin{F} & 1 & 110 & 1 \\ \hline
     \textlatin{0B} & \textlatin{Null} & 0 & 000 & \textlatin{x} \\ \hline
     \cdots & \textlatin{Null} & 0 & 000 & \textlatin{x} \\ \hline
     \textlatin{BF} & \textlatin{Null} & 0 & 000 & \textlatin{x} \\ \hline
     \textlatin{C0} & 0 & 1 & 110 & 1 \\ \hline
     \textlatin{C1} & 2 & 1 & 100 & \textlatin{x} \\ \hline
     \textlatin{C2} & \textlatin{Null} & 0 & 110 & \textlatin{x} \\ \hline
     \textlatin{C3} & \textlatin{Null} & 0 & 110 & \textlatin{x} \\ \hline
     \textlatin{C4} & \textlatin{Null} & 0 & 000 & \textlatin{x} \\ \hline
     \cdots & \textlatin{Null} & 0 & 000 & \textlatin{x} \\ \hline
     \textlatin{FD} & \textlatin{Null} & 0 & 000 & \textlatin{x} \\ \hline
     \textlatin{FE} & \textlatin{D} & 1 & 110 & 1 \\ \hline
     \textlatin{FF} & 1 & 1 & 110 & 0 \\ \hline
\end{tabular}_{ \ 256 \times 5 \ \textlatin{translation table}}
\]

$ \ \ \ \ \ \ \ \ \ \ \ \ {* \ \textlatin{Valid Bit = 0} \ \rightarrow \ \textlatin{Protection Bits = 000} \ ...}$

\vspace{1cm}


\subsubsection*{\textbf{γ)}}
\subsubsection*{\textbf{Έχουμε:}} 

\vspace{0.5cm}



\hspace{0.385cm} \text{- Εικονική διεύθυνση:} \ \ \textlatin{01038 (fetch)} \ $\implies$ \ \text{Φυσική διεύθυνση μετάφραση:} \ \ 4038.

\text{- Εικονική διεύθυνση:} \ \ \textlatin{0B0F4 (read)} \ $\implies$ \ \text{Τερματισμός προγράμματος - προσπάθεια προσπέλασης παράνομης διεύθ.}

\text{- Εικονική διεύθυνση:} \ \ \textlatin{C001C (write)} \ $\implies$ \ \text{Φυσική διεύθυνση μετάφραση:} \ \ \textlatin{001C}.

\text{- Εικονική διεύθυνση:} \ \ \textlatin{0292C (fetch)} \ $\implies$ \ \text{Επανεκίνηση προγράμματος - η σελίδα βρίσκεται στον δίσκο.}

\text{- Εικονική διεύθυνση:} \ \ \textlatin{00000 (read)} \ $\implies$ \ \text{Τερματισμός προγράμματος - προσπάθεια προσπέλασης παράνομης διεύθ.}

\text{- Εικονική διεύθυνση:} \ \ \textlatin{99F88 (read)} \ $\implies$ \ \text{Τερματισμός προγράμματος - προσπάθεια προσπέλασης παράνομης διεύθ.}

\text{- Εικονική διεύθυνση:} \ \ \textlatin{FE5D8 (write)} \ $\implies$ \ \text{Φυσική διεύθυνση μετάφραση:} \ \ \textlatin{D5D8}.

\text{- Εικονική διεύθυνση:} \ \ \textlatin{FF100 (fetch)} \ $\implies$ \ \text{Τερματισμός προγράμματος - προσπάθεια εκτέλεσης χωρίς πρόσβαση.}

\text{- Εικονική διεύθυνση:} \ \ \textlatin{C20CC (write)} \ $\implies$ \ \text{Επανεκίνηση προγράμματος - η σελίδα βρίσκεται στον δίσκο.}

\text{- Εικονική διεύθυνση:} \ \ \textlatin{CD0CC (write)} \ $\implies$ \ \text{Τερματισμός προγράμματος - προσπάθεια προσπέλασης παράνομης διεύθ.}

\text{- Εικονική διεύθυνση:} \ \ \textlatin{C0444 (read)} \ $\implies$ \ \text{Φυσική διεύθυνση μετάφραση:} \ \ \textlatin{0444}.

\text{- Εικονική διεύθυνση:} \ \ \textlatin{01FF4 (fetch)} \ $\implies$ \ \text{Φυσική διεύθυνση μετάφραση:} \ \ \textlatin{4FF4}.

\text{- Εικονική διεύθυνση:} \ \ \textlatin{C1FFC (write)} \ $\implies$ \ \text{Τερματισμός προγράμματος - προσπάθεια εκτέλεσης χωρίς πρόσβαση.}

\text{- Εικονική διεύθυνση:} \ \ \textlatin{008E4 (write)} \ $\implies$ \ \text{Τερματισμός προγράμματος - προσπάθεια προσπέλασης παράνομης διεύθ.}

\text{- Εικονική διεύθυνση:} \ \ \textlatin{C7700 (read)} \ $\implies$ \ \text{Τερματισμός προγράμματος - προσπάθεια προσπέλασης παράνομης διεύθ.}

\text{- Εικονική διεύθυνση:} \ \ \textlatin{01E40 (write)} \ $\implies$ \ \text{Τερματισμός προγράμματος - προσπάθεια εκτέλεσης χωρίς πρόσβαση.}








% -------------------------------------------------------------- %




 

\newpage

\text{   }
\vspace{-1.25cm}
\begin{center}
    {\underline{\Large \textbf{Άσκηση 12.5}}}
\end{center}
\vspace{0.345cm}



\subsubsection*{\textbf{αβγ)}}
\text{- Θέλουμε δύο επίπεδα, για κάθε διεργασία ένα πίνακα δευτέρου επιπέδου με 16 θέσεις. Έτσι, πρώτο επίπεδο: 256 / 16 = 16.}

\vspace{0.25cm}

\begin{center}
    \[
        \boxed{ \ \textlatin{Register: \ $*$ptr \ (pointer to the beginning of the level 1 translation table)} \ } 
    \]
\end{center}

\vspace{0.25cm}

\[
\text{Πίνακας πρώτου επιπέδου (\textlatin{L1}) 16 θέσεων (δεκαεξαδικό):}
\]
\[
\begin{aligned}
&\smash{\raisebox{3.75em}{\textlatin{$*$ptr \ \rightarrow \ \ \ }}}
\begin{tabular}{|c|c|c|}
    \hline
     \textbf{\textlatin{Index (4MS Virtual Address Bits)}} & \textbf{\textlatin{Valid Bit}} & \textbf{\textlatin{Pointer (To Level 2)}} \\ \hline
     0 & 1 & $*$\textlatin{ptr0} \\ \hline
     1 & 0 & \textlatin{x} \\ \hline
     \cdots & \cdots & \cdots \\ \hline
     \textlatin{B} & 0 & \textlatin{x} \\ \hline
     \textlatin{C} & 1 & $*$\textlatin{ptrC} \\ \hline
     \textlatin{D} & 0 & \textlatin{x} \\ \hline
     \textlatin{E} & 0 & \textlatin{x} \\ \hline
     \textlatin{F} & 1 & $*$\textlatin{ptrF} \\ \hline
\end{tabular}_{\ \textlatin{L1} \ (16 \times 3)}
\end{aligned}
\]

\[
\text{Πίνακες δευτέρου επιπέδου (\textlatin{L2}) 16 θέσεων (δεκαεξαδικό):}
\]
\[
\begin{aligned}
&\smash{\raisebox{4.5em}{\textlatin{$*$ptr0 \ \rightarrow \ \ \ }}}
\begin{tabular}{|c|c|c|c|c|c|}
    \hline
     \textbf{\textlatin{Index (4 VA LSB)}} & \textbf{\textlatin{Physical Page}} & \textbf{\textlatin{Valid B}} & \textbf{\textlatin{Protection B}} & \textbf{\textlatin{Dirty B}} & \textbf{\textlatin{Reference B}}  \\ \hline
     0 & \textlatin{Null} & 0 & 000 & \textlatin{x} & \textlatin{x}  \\ \hline 
     1 & 4 & 1 & 101 & \textlatin{x} & \textlatin{x}  \\ \hline
     2 & \textlatin{Null} & 0 & 101 & \textlatin{x} & \textlatin{x} \\ \hline
     \cdots & \textlatin{Null} & 0 & 101 & \textlatin{x} & \textlatin{x} \\ \hline
     9 & \textlatin{Null} & 0 & 101 & \textlatin{x} & \textlatin{x} \\ \hline
     \textlatin{A} & \textlatin{F} & 1 & 110 & 1 & \textlatin{x} \\ \hline
     \textlatin{B} & \textlatin{Null} & 0 & 000 & \textlatin{x} & \textlatin{x} \\ \hline
     \cdots & \textlatin{Null} & 0 & 000 & \textlatin{x} & \textlatin{x} \\ \hline
     \textlatin{F} & \textlatin{Null} & 0 & 000 & \textlatin{x} & \textlatin{x} \\ \hline
\end{tabular}_{\ \textlatin{L2}_0 \ (16 \times 6)}
\end{aligned}
\]
\[
\begin{aligned}
&\smash{\raisebox{3.1em}{\textlatin{$*$ptrC \ \rightarrow \ \ \ }}}
\begin{tabular}{|c|c|c|c|c|c|}
    \hline
     \textbf{\textlatin{Index (4 VA LSB)}} & \textbf{\textlatin{Physical Page}} & \textbf{\textlatin{Valid B}} & \textbf{\textlatin{Protection B}} & \textbf{\textlatin{Dirty B}} & \textbf{\textlatin{Reference B}}  \\ \hline
     \textlatin{0} & 0 & 1 & 110 & 1 & \textlatin{x} \\ \hline
     \textlatin{1} & 2 & 1 & 100 & \textlatin{x} & \textlatin{x} \\ \hline
     \textlatin{2} & \textlatin{Null} & 0 & 110 & \textlatin{x} & \textlatin{x} \\ \hline
     \textlatin{3} & \textlatin{Null} & 0 & 110 & \textlatin{x} & \textlatin{x} \\ \hline
     \textlatin{4} & \textlatin{Null} & 0 & 000 & \textlatin{x} & \textlatin{x} \\ \hline
     \cdots & \textlatin{Null} & 0 & 000 & \textlatin{x} & \textlatin{x} \\ \hline
     \textlatin{F} & \textlatin{Null} & 0 & 000 & \textlatin{x} & \textlatin{x} \\ \hline
\end{tabular}_{\ \textlatin{L2}_\textlatin{C} \ (16 \times 6) \ \ \ \ \ }
\end{aligned}
\]
\[
\begin{aligned}
&\smash{\raisebox{1.999em}{\textlatin{$*$ptrF \ \rightarrow \ \ \ }}}
\begin{tabular}{|c|c|c|c|c|c|}
    \hline
     \textbf{\textlatin{Index (4 VA LSB)}} & \textbf{\textlatin{Physical Page}} & \textbf{\textlatin{Valid B}} & \textbf{\textlatin{Protection B}} & \textbf{\textlatin{Dirty B}} & \textbf{\textlatin{Reference B}}  \\ \hline
    \textlatin{0} & \textlatin{Null} & 0 & 000 & \textlatin{x} & \textlatin{x} \\ \hline
     \cdots & \textlatin{Null} & 0 & 000 & \textlatin{x} & \textlatin{x} \\ \hline
     \textlatin{D} & \textlatin{Null} & 0 & 000 & \textlatin{x} & \textlatin{x} \\ \hline
     \textlatin{E} & \textlatin{D} & 1 & 110 & 1 & \textlatin{x} \\ \hline
     \textlatin{F} & 1 & 1 & 110 & 0 & \textlatin{x} \\ \hline
\end{tabular}_{\ \textlatin{L2}_\textlatin{F} \ (16 \times 6) \ \ }
\end{aligned}
\]


\vspace{0.5cm}


\subsubsection*{\textbf{δ)}}

\hspace{0.35cm} \text{Έχουμε:} \ 
$ \textlatin{Total bits now} = \textlatin{L1} + \textlatin{L2} + \textlatin{Pages} = (16)(4+1+4) + (3)(16)(4+4+1+3+1+1) + (6)(4\textlatin{KB)} = 197.424. $

\vspace{0.01cm}

\text{Έχουμε:} \ 
$ \textlatin{Total bits then} = \textlatin{Table} + \textlatin{Pages} = (256)(8+4+1+3+1+1) + (6)(4\textlatin{KB)} = 201.216. $ \hspace{10cm}

\vspace{0.01cm}

\text{Τελικά, αφού} \ 197.424 \( < \) 201.216 \ υπάρχει οικονομία στο χώρο μνήμης που καταλαμβάνεται ...





% -------------------------------------------------------------- %




 

\newpage

\text{   }
\vspace{-1.25cm}
\begin{center}
    {\underline{\Large \textbf{Άσκηση 12.6}}}
\end{center}
\vspace{0.25cm}


\subsubsection*{\textbf{Δεδομένα:}}
\begin{itemize}
    \item Εικονικές διευθύνσεις: 64 \textlatin{bits} αλλά χρησιμοποιούνται τα 48 \textlatin{LS Bits}.
    \item Χώρος εικονικών διευθύνσεων: 256 \textlatin{TeraBytes} ανά διεργασία.
    \item Μέγεθος σελίδας: 4\textlatin{KB} ($2^{12}$ \textlatin{bytes}) άρα 12 \textlatin{LS bits} κάθε εικονικής διεύθυνσης για \textlatin{offset}.
    \item Τρία (3) επίπεδα πινάκων μετάφρασης ανά διεργασία.
    \item Μέγεθος/Θέσεις πίνακα μετάφρασης: 4096 (άρα κοιτάζει 12 από τα \textlatin{bits} του αριθμού εικονικής σελίδας).
    \item Δυνατότητα επέκτασης έως 256 \textlatin{GBytes} φυσικής μνήμης ...
    \item (38 \textlatin{bits}
    φυσική διεύθυνση ή 64Μ φυσικές σελίδες (26 \textlatin{bits} αριθμός φυσικής
    σελίδας) επί 4 \textlatin{KBytes} ανά σελίδα)
\end{itemize}

\subsubsection*{\textbf{Έχουμε:}}

\text{Έχουμε 64 \textlatin{bits} και χρησιμοποιούμε 48 - 12 (\textlatin{LSBits}) = 36 / 3 (\textlatin{Levels}) = 12 \textlatin{bits} κάθε επίπεδο / 4 = 3 \textlatin{bits} το κάθε \textlatin{Index} ...}

\vspace{-0.35cm}
\begin{center}
    \[
        \boxed{ \ \textlatin{Virtual Address: 16 Unused Bits \& 12 L1 Bits \& 12 L2 Bits \& 12 L3 Bits \& 12 Not Translated Bits = 64 Bits} \ } 
    \]
\end{center}
\vspace{-1cm}
\begin{center}
    \[
        \boxed{ \ \textlatin{Register ($*$pointer pointing to the beginning of the level 1 translation table)} \ } 
    \]
\end{center} 
\[
\begin{aligned}
&\smash{\raisebox{1.75em}{\textlatin{\ $*$ptr \ \rightarrow \ \ \ }}}
\begin{tabular}{|c|c|c|}
    \hline
     \textbf{\textlatin{Index}} & \textbf{\textlatin{Valid}} & \textbf{\textlatin{Pointers}} \\ \hline
     000 & 1 & $*$\textlatin{ptr0} \\ \hline
     001 & 0 & \textlatin{Null} \\ \hline
     \cdots & \cdots & \cdots \\ \hline
     \textlatin{FFE} & 0 & \textlatin{Null} \\ \hline
     \textlatin{FFF} & 1 & $*$\textlatin{ptrF} \\ \hline
\end{tabular}_{\ \textlatin{L1} \ (4096 \times 3) \ \ }
\end{aligned}
\]
\[
\begin{aligned}
&\begin{array}{r}
\smash{\raisebox{0.25ex}{$*$\textlatin{ptr0} $\rightarrow$}} \\
\vspace{3ex} \\
\vspace{3ex} \\
\smash{\raisebox{2.75ex}{$*$\textlatin{ptrF} $\rightarrow$}}
\end{array}
\quad
\begin{tabular}{|c|c|c|}
    \hline
     \textbf{\textlatin{Index}} & \textbf{\textlatin{Valid}} & \textbf{\textlatin{Pointers}} \\ \hline
     000 & 1 & $*$\textlatin{ptr00} \\ \hline
     \cdots & \cdots & \cdots \\ \hline
     \textlatin{FFF} & 1 & $*$\textlatin{ptr0F} \\ \hline
     \cdots & \cdots & \cdots \\ \hline
     000 & 1 & $*$\textlatin{ptrF0} \\ \hline
     \cdots & \cdots & \cdots \\ \hline
     \textlatin{FFF} & 1 & $*$\textlatin{ptrFF} \\ \hline
\end{tabular}_{\ \textlatin{L2} \ ((2^{12} \cdot 2^{12}) \times 3)}
\end{aligned}
\]
\[
\begin{aligned}
&\begin{array}{r}
\smash{\raisebox{3ex}{$*$\textlatin{ptr00} $\rightarrow$}} \\
\vspace{5ex} \\
\smash{\raisebox{1.5ex}{$*$\textlatin{ptr0F} $\rightarrow$}} \\
\vspace{5ex} \\
\smash{\raisebox{-0.1ex}{$*$\textlatin{ptrF0} $\rightarrow$}} \\
\vspace{5ex} \\
\smash{\raisebox{-0.5ex}{$*$\textlatin{ptrFF} $\rightarrow$}}
\end{array}
\quad
\begin{tabular}{|c|c|c|}
    \hline
     \textbf{\textlatin{Index}} & \textbf{\textlatin{All Bits}} & \textbf{\textlatin{Page}} \\ \hline
     000 & \textlatin{xxxxxx} & \textlatin{Null} \\ \hline
     \cdots & \cdots & \cdots \\ \hline
     \textlatin{FFF} & \textlatin{xxxxxx} & \textlatin{-x-} \\ \hline
     \cdots & \cdots & \cdots \\ \hline
     000 & \textlatin{xxxxxx} & \textlatin{Null} \\ \hline
     \cdots & \cdots & \cdots \\ \hline
     \textlatin{FFF} & \textlatin{xxxxxx} & \textlatin{-x-} \\ \hline
     \cdots & \cdots & \cdots \\ \hline
     000 & \textlatin{xxxxxx} & \textlatin{Null} \\ \hline
     \cdots & \cdots & \cdots \\ \hline
     \textlatin{FFF} & \textlatin{xxxxxx} & \textlatin{-x-} \\ \hline
     \cdots & \cdots & \cdots \\ \hline
     000 & \textlatin{xxxxxx} & \textlatin{Null} \\ \hline
     \cdots & \cdots & \cdots \\ \hline
     \textlatin{FFF} & \textlatin{xxxxxx} & \textlatin{-x-} \\ \hline
\end{tabular}_{\ \textlatin{L3} \ (2^{36} \times 3) \quad \quad \ \ \ }
\end{aligned}
\]
\begin{center}
    \[
        \boxed{ \ \textlatin{Physical Address: 26 L3 Page Bits \& 12 Virtual Adrress Bits = 48 Total Used Bits} \ } 
    \]
\end{center} 





 

% -------------------------------------------------------------- %




 

\newpage

\text{   }
\vspace{-1.25cm}
\begin{center}
    {\underline{\Large \textbf{Άσκηση 12.7}}}
\end{center}
\vspace{0.75cm}


\subsubsection*{\textbf{\alpha)}}


\[
\renewcommand{\arraystretch}{1.3}
\begin{tabular}{|c|c|}
    \hline
    \multicolumn{2}{|c|}{\textbf{\large \textlatin{TLB}}} \\ \hline
    \textbf{Πεδία (ανά θέση)} & \textbf{Μέγεθος} \\ \hline
    \textlatin{PID} & \textlatin{8 bits} \\ \hline
    \textlatin{Virtual Address} & \textlatin{8 bits} \\ \hline
    \textlatin{Physical Address} & \textlatin{4 bits} \\ \hline
    \textlatin{Valid Bit} & \textlatin{1 bit} \\ \hline
    \textlatin{Protection Bits} & \textlatin{3 bits} \\ \hline
    \textlatin{Dirty Bit} & \textlatin{1 bit} \\ \hline
    \textlatin{Reference Bit} & \textlatin{1 bit} \\ \hline
\end{tabular}
\]


\vspace{1cm}

\subsubsection*{\textbf{\beta)}}


\[
\renewcommand{\arraystretch}{1.1}
\begin{tabular}{|c|c|c|c|c|c|c|}
\hline
\multicolumn{7}{|c|}{\textbf{\textlatin{TLB}}} \\ \hline
\textlatin{PID} & \textlatin{V. Address} & \textlatin{P. Address} & \textlatin{Valid B.} & \textlatin{Protection B.} & \textlatin{Dirty B.} & \textlatin{Reference B.} \\
\hline
\textlatin{3B} & 03 & 1 & 1 & 100 & 0 & 1 \\ % Κώδικας (read+exec)
\textlatin{B4} & 03 & Α & 1 & 110 & 1 & 1 \\ % Data r/w
\textlatin{3C} & 03 & 1 & 1 & 100 & 0 & 1 \\ % Κοινός κώδικας με 3B
\textlatin{3B} & \textlatin{FF} & 6 & 1 & 110 & 1 & 1 \\ % Στοίβα 3B
\textlatin{3C} & \textlatin{FF} & 2 & 1 & 110 & 1 & 1 \\ % Στοίβα 3C
\textlatin{A2} & \textlatin{C2} & \textlatin{F} & 1 & 010 & 1 & 1 \\ % Κοινόχρηστη μνήμη (write only)
\textlatin{A3} & \textlatin{E3} & \textlatin{F} & 1 & 100 & 0 & 1 \\ % Κοινόχρηστη μνήμη (read only)
\hline
\end{tabular}_{\ 7 \times 7}
\]



\vspace{1cm}

\subsubsection*{\textbf{\gamma)}}


\vspace{0.12345cm}
\noindent
\text{Οι διεργασίες \textlatin{3B} και \textlatin{3C} είναι προστατευμένες μεταξύ τους ως προς τα δεδομένα τους, επειδή κάθε μία χρησιμοποιεί διαφορε-} \\
\text{τικές φυσικές σελίδες για τα δεδομένα της (\textlatin{physical addresses} 3 και 4 αντίστοιχα). Έτσι, η \textlatin{3B} δεν μπορεί να} \\ 
\text{διαβάσει ή να τροποποιήσει τα δεδομένα της \textlatin{3C} και το αντίστροφο.} \\
\text{Ο κώδικας των δύο διεργασιών είναι κοινός και βρίσκεται σε μία κοινή φυσική σελίδα (\textlatin{physical address} 1), η οποία έχει προ-} \\
\text{στασία \textlatin{read-only} (\textlatin{protection bits} 100). Αυτό σημαίνει ότι ο κώδικας δεν μπορεί να τροποποιηθεί από καμία από τις διεργασίες,} \\ 
\text{αποτρέποντας αλλοιώσεις ή καταστροφές.} \\ 
\text{Η προστασία αυτή επιτυχάνεται λώγο του \textlatin{TLB}, το οποίο ελέγχει την πρόσβαση βάσει των} \\
\text{\textlatin{protection bits} και των \textlatin{physical addresses}. Με αυτόν τον τρόπο, ενώ οι διεργασίες μοιράζονται τον ίδιο κώδικα για οικονομία} \\
\text{μνήμης, διατηρούν πλήρη ανεξαρτησία και ασφάλεια στα δεδομένα τους.} \\
\text{}






% ---------------------------------------------------------------------------------------

\newpage

\begin{textblock}{10}(14.401, 15.501)
    \boxed{\textbf{\LaTeX{}}}
\end{textblock}

\begin{center}
    \vspace*{\fill}
    \textbf{\LARGE{ΤΕΛΟΣ}} \\[0.001cm]
    \textbf{\hspace{0.1cm}(ευχαριστώ)}
    \\[0.1cm]
    \vspace*{\fill}
\end{center}

% ---------------------------------------------------------------------------------------

\end{document}

% -------------------------------------------------------------- %
%                           The End :)                           %
% -------------------------------------------------------------- %
